\documentclass[14pt,a4paper,oneside]{extreport}
\usepackage{gost732}
% All names and data in this example are educational and fictional.
\newcommand\ReportMinistry{Учебный пример оформления}
\newcommand\ReportOrganization{Условная исследовательская организация}
\newcommand\ReportShortOrganization{УИО}
\newcommand\ReportUDC{не определялся}
\newcommand\ReportResearchID{не присваивался}
\newcommand\ReportID{не присваивался}
\newcommand\ReportTitle{Исследование точности измерения времени}
\newcommand\ReportType{заключительный}
\newcommand\ReportProgram{Учебный пример}
\newcommand\ReportApproverRole{Руководитель организации}
\newcommand\ReportApprover{И.О. Фамилия}
\newcommand\ReportLeaderRole{Руководитель НИР}
\newcommand\ReportLeader{И.О. Фамилия}
\newcommand\ReportAuthorRole{Исполнитель}
\newcommand\ReportAuthor{И.О. Фамилия}
\newcommand\ReportCity{Москва}
\newcommand\ReportYear{2026}

\begin{document}
\begin{GostTitlePage}
\begin{center}
\ReportMinistry\par
\MakeUppercase{\ReportOrganization}\par
(\MakeUppercase{\ReportShortOrganization})\par
\end{center}
\vspace{2\baselineskip}
\noindent УДК \ReportUDC\\
Рег. № НИР \ReportResearchID\\
Рег. № отчёта \ReportID\par
\vspace{2\baselineskip}
\noindent\hfill\begin{minipage}[t]{.47\textwidth}
УТВЕРЖДАЮ\\
\ReportApproverRole\\
\rule{17mm}{.4pt}\ \ReportApprover\\
\rule{25mm}{.4pt}
\end{minipage}\par
\vspace{2\baselineskip}
\begin{center}
ОТЧЕТ\\
О НАУЧНО-ИССЛЕДОВАТЕЛЬСКОЙ РАБОТЕ\par
\MakeUppercase{\ReportTitle}\par
(\ReportType)\par
\ReportProgram\par
\end{center}
\vfill
\noindent\begin{tabular}{@{}p{.42\textwidth}p{.13\textwidth}p{.37\textwidth}@{}}
\ReportLeaderRole & \rule{14mm}{.4pt} & \ReportLeader\\[1\baselineskip]
\ReportAuthorRole & \rule{14mm}{.4pt} & \ReportAuthor
\end{tabular}\par
\vfill
\begin{center}\ReportCity\space\ReportYear\end{center}
\end{GostTitlePage}


% One author: omit a separate list under 5.2.2; provide their details on title.
\GostStructural{РЕФЕРАТ}
Отчёт 8 с., 1 кн., 1 рис., 2 табл., 1 источник, 1 прил.

\GostKeywords{измерение, время, среднее значение, погрешность, повторяемость}

Объект исследования — результаты повторного измерения времени.

Цель работы — оценить среднее значение измеряемой величины.

В работе использован метод повторных наблюдений. Получен пример расчёта среднего значения. Приведённые демонстрационные данные следует заменить фактическими результатами исследования.

\GostTOC

\GostStructural{ВВЕДЕНИЕ}
Рассматривается определение среднего значения по результатам повторных измерений. Исходные данные представлены в учебном наборе \GostCite{src:units}.

\chapter{Метод исследования}
\section{Исходные данные}
Результаты представлены в таблице~\ref{tab:data}. Схема обработки приведена на рисунке~\ref{fig:process}. Дополнительные сведения находятся в приложении~\ref{app:data}.

\begin{table}[!htbp]
\caption{Результат измерения}\label{tab:data}
\centering
\begin{tabular}{|>{\centering\arraybackslash}p{60mm}|>{\centering\arraybackslash}p{60mm}|}
\hline
Номер наблюдения & Время, с\\\hline
1 & 10,1\\\hline
2 & 10,3\\\hline
\end{tabular}
\end{table}

\begin{figure}[!htbp]
\centering
\setlength\unitlength{1mm}
\begin{picture}(110,22)
\put(0,4){\framebox(40,14){Наблюдения}}
\put(40,11){\vector(1,0){25}}
\put(65,4){\framebox(45,14){Среднее значение}}
\end{picture}
\caption{Последовательность обработки}\label{fig:process}
\end{figure}

\section{Расчёт}
Среднее значение определяется по формуле~\eqref{eq:mean}.
\begin{equation}\label{eq:mean}
\overline{t}=\frac{1}{n}\sum_{i=1}^{n}t_i.
\end{equation}
\GostWhere{$\overline{t}$ — среднее время, с;\\
$n$ — число наблюдений;\\
$t_i$ — время отдельного наблюдения, с.}

\GostNote{Вычисление следует повторить для фактических исходных данных.}

\GostStructural{ЗАКЛЮЧЕНИЕ}
Для двух демонстрационных наблюдений среднее время составляет 10,2~с. Выводы об эффективности и применимости требуют действительных результатов исследования.

\begin{GostSources}
\GostSource{src:units}{Учебные данные повторных наблюдений: локальный набор. — 2026. — 2 наблюдения.}
\end{GostSources}

\GostAppendix{Дополнительные данные}{справочное}{app:data}
В таблице~\ref{tab:extra} показаны дополнительные наблюдения.
\begin{table}[!htbp]
\caption{Дополнительное наблюдение}\label{tab:extra}
\centering
\begin{tabular}{|p{60mm}|>{\centering\arraybackslash}p{60mm}|}
\hline
Наблюдение & Время, с\\\hline
Контрольное & 10,2\\\hline
\end{tabular}
\end{table}
\end{document}
